\subsection{Exercises 2-2}

\num What is the impact of memory coalescing on DRAM access efficiency in CUDA?

\begin{tcolorbox}
	When all threads of a warp execute a load instruction, if all accessed locations fall into the same burst section (coalesced access), only one DRAM request will be made and the access is fully coalesced.
\end{tcolorbox}

\num What is the purpose of using shared memory in tiled matrix multiplication? Is the corner turning technique relevant for this purpose?

\begin{tcolorbox}
	Reading A and B in a matrix multiplication results in A read in non-coalesced way and B read in a coalesced way. To improve performance, we copy the data into shared memory, allowing both input matrices to be accessed coalesced in both cases (corner turning).
\end{tcolorbox}

\num Please describe how priority scatter works.

\begin{tcolorbox}
	Given a collection of input data, and a collection of location, \textbf{scatter} copies the input data in the output data, at the location specified by the locations collection.

	The location array could map two or more input elements to the same location. In this case, priority scatter solves this write conflict by allowing only the processor with highest priority to write.
\end{tcolorbox}

\num How can a CUDA programmer judge if an access is coalesced?

\begin{tcolorbox}
	By checking if consecutive threads in a warp are accessing consecutive memory addresses. Mathematically, the access is coalesced if the memory accesses are in the form of \texttt{Mem[Base + ThreadIdx.x]}. This creates a sequential pattern that allows the memory controller to merge the individual requests into a single transaction.

	If the threads access memory e.g. randomly, the hardware must split the operation into multiple separate transactions, signaling uncoalesced access.
\end{tcolorbox}

\num Describe the corner turning technique. How the size of the shared memory is defined in this technique?

\begin{tcolorbox}
	\textit{See corner turning technique definition in a previous question}

	The allocated size is primarily defined by the tile dimensions (e.g., $TILE\_WIDTH \times TILE\_WIDTH$). However, in \textbf{Corner Turning}, the size is frequently defined as
	$$
		TILE\_WIDTH \times (TILE\_WIDTH + 1)
	$$
	This extra column of padding shifts the data layout in shared memory to ensure that when threads read down a column, they access different memory banks, preventing Bank Conflicts.
\end{tcolorbox}

\num Does it make sense to use the corner-turning optimization in a non-GPU-based architecture? Why?

\begin{tcolorbox}
	\textbf{Yes}. The corner-turning technique is effective also on CPUs to maximize spatial locality in the cache. CPUs fetch memory in cache lines; if a program accesses data column-wise in a row-major array, it loads an entire cache line but uses only a single value, wasting bandwidth and causing cache pollution.

	Corner turning reorganizes the data so that the CPU reads contiguous memory addresses, enabling the use of SIMD instructions, which require aligned data for maximum performance.
\end{tcolorbox}

\num Please briefly describe the differences between Memory coherence and memory consistency.

\begin{tcolorbox}
	\begin{itemize}
		\item Memory \textbf{coherence} ensures that \textbf{all processors see the most recent update} to that specific location (everyone agrees that a specific variable has a specific value).
		\item Memory \textbf{consistency} defines the rules for the \textbf{ordering of operations} across multiple distinct memory addresses, ensureing the timeline of events makes sense across the entire program.
	\end{itemize}
\end{tcolorbox}

\num Please describe how merge scatter works.

\begin{tcolorbox}
	The scatter pattern, given a collection of input data and a collection of write locations, writes the data to the output collection at the indices specified by the location array.

	\textbf{Merge} is used to handle collisions when multiple input elements attempt to write to the same output location. It combines the conflicting writes using an operator (e.g., addition: the final value is the sum of the conflicting input values). This operator must be \textbf{associative} and \textbf{commutative}.
\end{tcolorbox}

\num Is DRAM burst affecting the corner-turning optimization?

\begin{tcolorbox}
	\textbf{Yes}. DRAM burst necessitates corner-turning optimization, because DRAM transfers data in fixed-size chunks (bursts), typically filling a cache line in one transaction. Accessing memory in a strided pattern like reading column-wise, pulls in a full burst but consumes only a single useful value, wasting the rest of the available bandwidth.

	Corner-turning reorganizes the data into a contiguous layout, ensuring that when a DRAM burst is fetched, every byte within it is utilized by the processor, maximizing memory throughput.
\end{tcolorbox}

\num Please briefly describe what DRAM Banking is meant for. How can CUDA take advantage of DRAM banking?


\begin{tcolorbox}
	DRAM banking partitions physical memory into independent sub-arrays (banks) that can operate in parallel, allowing the memory controller to issue a command to another bank while one is busy (e.g., activating a row).

	CUDA takes advantage of this because when a warp performs a coalesced memory access (reading contiguous data), the requests are naturally distributed across multiple active banks, preventing serialization (bank conflicts) and saturating the global memory bus.
\end{tcolorbox}

\num Describe briefly the shared memory programming model with threads and two of its implementations.

\begin{tcolorbox}
	The \textbf{shared memory} programming model involves multiple concurrent threads operating within a single unified address space, where they communicate implicitly by reading/writing to the same variables. This model relieson synchronization primitives (like locks or atomic operations) to prevent race conditions.

	Two key implementations are
	\begin{itemize}
		\item POSIX Threads (Pthreads): low-level, explicit API standard used in C/C++ that gives developers fine-grained control over thread lifecycle and synchronization;
		\item OpenMP: high-level standard that automates thread creation and work distribution via compiler pragmas (e.g., \texttt{\#pragma omp parallel})
	\end{itemize}
\end{tcolorbox}

\num Describe in detail the effect of the following OpenMP pragma.

\begin{minted}[autogobble]{text}
    #pragma omp parallel for shared(tmp) private(i, j) if (!embedded)
\end{minted}

\begin{tcolorbox}
	Parallelized the immediately following \texttt{for} loop:

	\begin{enumerate}
		\item \texttt{\#pragma omp parallel for}: creates a team of threads and divides the iterations of the following \texttt{for} loop amond those threads. At the end of the loop, there is an implicit \textbf{barrier} where all threads wait for each other to finish before the master thread resumes execution.
		\item \texttt{shared(tmp)}: means that the variable \texttt{tmp} exists in a single memory location shared by all threads.
		\item \texttt{private(i, j)}: means that each thread creates its own local instance of variables \texttt{i} and \texttt{j}. \texttt{i} is implicitly handled by OpenMP, but \texttt{j} is \textbf{uninitialized}.
		\item \texttt{if (!embedded)}: determines at runtime whether to execute in parallel or sequentially, based on the value of the \texttt{embedded} boolean variable.
	\end{enumerate}
\end{tcolorbox}

\num Which OpenMP pragma allows offloading part of the computation to an accelerator? Which clause is used to manage data transfers?

\begin{tcolorbox}
	\begin{itemize}
		\item \texttt{\#pragma omp target} is used to offload a code region to an accelerator (e.g., GPU).
		\item To manage the data transfer between the host and the device, we use the \texttt{map} clause.
	\end{itemize}
\end{tcolorbox}

\num List the factors that determine the number of threads in an OpenMP parallel region.

\begin{tcolorbox}
	The factors listed in order from highest precedence to lowest:
	\begin{itemize}
		\item \texttt{num\_threads} clause in a pragma declaration;
		\item \texttt{omp\_set\_num\_threads()} function can be called by the program before the parallel region;
		\item \texttt{OMP\_NUM\_THREADS} environment variable set by the user in the OS shell before running the program;
		\item \textbf{system default} if none of the above are set.
	\end{itemize}
\end{tcolorbox}

\num Which errors will arise compiling the following code snippet?

\begin{minted}[autogobble]{text}
    #pragma omp parallel for shared(num_steps) reduction(+:sum) default(none)
    for (i=1;i<= num_steps; i++)
    {
        x = array[i] * 2;
        sum = sum + x;
    }
\end{minted}

\begin{tcolorbox}
	The compiler will issue error regarding the \textbf{scope of the data}. The clause \texttt{default(none)} enforces that every variable inside the parallel region must be explicitly categorized (e.g., \texttt{shared}, \texttt{private}), but \texttt{x} and \texttt{array} are missing from the directives.

	Note: \texttt{i} is implicitly private in a \texttt{parallel for} loop.
\end{tcolorbox}

\num Describe the way in which 15 iterations are divided among 10 threads with a static schedule and a chunk size of 2, and with a dynamic schedule with a chunk size of 1.

\begin{tcolorbox}
	\begin{itemize}
		\item \textbf{static schedule}: the iterations are divided into fixed blocks of 2 and assigned to threads in a \textbf{round robin} fashion: thread 0 gets iterations \texttt{0, 1}, thread 1 gets \texttt{2, 3}, etc. Note how threads 8 and 9 receive no work.
		\item \textbf{dynamic schedule}: the iterations are put in a queue and threads request 1 iteration at a time. Initially, all 10 threads (0–9) will likely pick up one iteration each. The remaining 5 iterations will be assigned one by one to whichever threads finish their first task fastest.
	\end{itemize}
\end{tcolorbox}

\num It is always necessary to place an explicit barrier at the end of an OpenMP work sharing construct. Is this statement true or false? Explain.

\begin{tcolorbox}
	\textbf{FALSE}. This statement is False.

	By default, OpenMP automatically inserts an implicit barrier at the end of most work-sharing constructs (like \texttt{for}, \texttt{sections}, etc) to ensure all threads finish their assigned tasks before proceeding.
\end{tcolorbox}

\num Describe what is the effect of the following OpenMP directive:

\begin{minted}[autogobble]{text}
    #pragma omp for simd schedule(simd:static, 3)
\end{minted}

\begin{tcolorbox}
	It splits the loop iterations among the available threads (the \texttt{for} part) and simultaneously instructs the compiler to execute the iterations within each thread using SIMD vector instructions (the \texttt{simd} part).

	The clause \texttt{schedule(simd:static, 3)} dictates that the iterations are divided into chunks of exactly 3, which are assigned to threads in a \textbf{static round-robin} fashion (e.g., Thread 0 gets iterations 0-2, Thread 1 gets 3-5, etc.). The \texttt{simd} modifier in the schedule ensures the chunk boundaries are handled consistently with SIMD alignment.
\end{tcolorbox}
